% Naturality under refinement and Rayleigh-type monotonicity for Stokes energies.

\begin{proposition}[粗粒化细化下的 Stokes 自然性与 Rayleigh--Stokes 单调性]\label{prop:fold-coarsegraining-naturality-rayleigh-stokes}
\leanverified{paper\_fold\_coarsegraining\_naturality\_rayleigh\_stokes}
设 \(n\ge 1\) 且 \(\Omega_n=\{0,1\}^n\)。
给定两级满射粗粒化 \(\tilde f:\Omega_n\twoheadrightarrow \tilde X\) 与 \(f:\Omega_n\twoheadrightarrow X\)，并存在满射 \(\rho:\tilde X\twoheadrightarrow X\) 使 \(f=\rho\circ \tilde f\)。
记 \(\mathsf B_{\tilde f}\)、\(\mathsf B_f\) 为相应边界邻接多图（定义 \ref{def:hypercube-boundary-multigraph}）。
\begin{enumerate}
  \item \textbf{（坐标净通量的自然性）}
  令 \(F^{(i)}_{\tilde f}\in C_1(\mathsf B_{\tilde f};\ZZ)\) 与 \(F^{(i)}_{f}\in C_1(\mathsf B_{f};\ZZ)\) 为命题 \ref{prop:fold-hypercube-godel-hodge-torsor} 给出的第 \(i\) 坐标自然整数流。
  定义推前算子 \(\rho_\#:C_1(\mathsf B_{\tilde f};\ZZ)\to C_1(\mathsf B_f;\ZZ)\) 为：对任意有向边 \(u\to v\)（\(u,v\in X\)）取
  \[
  (\rho_\#\Phi)(u\to v):=\sum_{\substack{\tilde u\in\rho^{-1}(u)\\ \tilde v\in\rho^{-1}(v)}}\Phi(\tilde u\to \tilde v).
  \]
  则对每个 \(i\) 有 \(\rho_\#F^{(i)}_{\tilde f}=F^{(i)}_{f}\)，并且边界算子满足交换律 \(\partial(\rho_\#\Phi)=\rho_\#(\partial\Phi)\)。
  \item \textbf{（Rayleigh--Stokes 单调性）}
  在 \(\mathsf B_{\tilde f}\) 的每条边 \(\tilde e\) 上取导纳 \(\tilde c_{\tilde e}>0\)，并对粗图边 \(e\) 定义推前导纳
  \[
  c_e:=\sum_{\tilde e\mapsto e}\tilde c_{\tilde e}.
  \]
  在 \(C_1(\mathsf B_f;\RR^n)\) 上定义能量
  \[
  \mathcal E_c(g):=\sum_{e\in E(\mathsf B_f)}\frac{\|g(e)\|^2}{c_e},
  \]
  并对守恒散度数据 \(b\in C_0(\mathsf B_f;\RR^n)\) 定义最小能量
  \[
  \mathfrak E_f(b):=\min\{\mathcal E_c(g):\ g\in C_1(\mathsf B_f;\RR^n),\ \partial g=b\}.
  \]
  在细图上同理定义 \(\mathfrak E_{\tilde f}(\tilde b)\)。
  令细化分配集合
  \[
  \mathcal B_{\tilde f}(b):=\{\tilde b\in C_0(\mathsf B_{\tilde f};\RR^n):\ \rho_\#\tilde b=b\}.
  \]
  则成立不等式
  \[
  \inf_{\tilde b\in\mathcal B_{\tilde f}(b)}\ \mathfrak E_{\tilde f}(\tilde b)\ \ge\ \mathfrak E_f(b).
  \]
\end{enumerate}
\end{proposition}
\begin{proof}
\textbf{（自然性）}
由定义，粗纤维按细纤维分解；超立方体第 \(i\) 坐标割边推前到边界邻接多图时，跨越粗纤维的割边是跨越相应细纤维割边的不交并，且取向保持。
因此坐标整数流在聚合下逐边相加，得到 \(\rho_\#F^{(i)}_{\tilde f}=F^{(i)}_{f}\)。
边界算子交换律为定义的直接重排。
\par
\textbf{（单调性）}
对任意 \(\tilde b\in\mathcal B_{\tilde f}(b)\)，取任意 \(\tilde g\in C_1(\mathsf B_{\tilde f};\RR^n)\) 满足 \(\partial\tilde g=\tilde b\)，并令 \(g:=\rho_\#\tilde g\)。
由 \(\partial(\rho_\#\cdot)=\rho_\#(\partial\cdot)\) 得 \(\partial g=b\)。
对每条粗边 \(e\) 记提升边集 \(\{\tilde e:\tilde e\mapsto e\}\)，则向量值 Cauchy--Schwarz 给出
\[
\frac{\|g(e)\|^2}{c_e}
=\frac{\bigl\|\sum_{\tilde e\mapsto e}\tilde g(\tilde e)\bigr\|^2}{\sum_{\tilde e\mapsto e}\tilde c_{\tilde e}}
\le
\sum_{\tilde e\mapsto e}\frac{\|\tilde g(\tilde e)\|^2}{\tilde c_{\tilde e}}.
\]
对全部粗边求和得 \(\mathcal E_c(g)\le \mathcal E_{\tilde c}(\tilde g)\)，从而
\[
\mathfrak E_f(b)\le \mathcal E_c(g)\le \mathcal E_{\tilde c}(\tilde g).
\]
对满足 \(\partial\tilde g=\tilde b\) 的 \(\tilde g\) 取极小，再对 \(\tilde b\in\mathcal B_{\tilde f}(b)\) 取下确界，即得结论。证毕。
\end{proof}

\endinput
